% Generated draft: requires mathematical and source review.
\documentclass[11pt]{article}
\usepackage[margin=25mm]{geometry}
\usepackage{fontspec}
\setmainfont{Latin Modern Roman}
\usepackage{amsmath,amssymb,mathtools,unicode-math}
\setmathfont{Latin Modern Math}
\usepackage{enumitem,needspace,xurl,hyperref}
\setlength{\parindent}{0pt}
\setlength{\parskip}{5pt}
\setlength{\emergencystretch}{3em}
\newcommand{\N}{\mathbb{N}}
\newcommand{\Z}{\mathbb{Z}}
\newcommand{\Q}{\mathbb{Q}}
\newcommand{\R}{\mathbb{R}}
\newcommand{\C}{\mathbb{C}}
\newcommand{\F}{\mathbb{F}}
\newcommand{\A}{\mathbb{A}}
\newcommand{\PP}{\mathbb{P}}
\newcommand{\E}{\mathbb{E}}
\newcommand{\eps}{\varepsilon}
\newcommand{\dd}{\,\mathrm{d}}
\begin{document}
\section*{UnsolvedMath ID: 20002443}
\textbf{Source problem code:} catalogue entry
\subsection{Definitions and mathematical statement}
The recovered R. Pego item concerns the Toda lattice
\[
 \ddot q_k=e^{q_{k+1}-q_k}-e^{q_k-q_{k-1}},\qquad k\in\Z,
\]
and its one-soliton
\[
 q_k(t)=Q_\beta(k-ct),\qquad
 Q_\beta(z)=\log\frac{\cosh(\beta(z+1))}{\cosh(\beta z)},
 \qquad\beta>0,\quad c=\frac{\sinh\beta}{\beta}>1.
\]
The source asks for its spectral stability properties. The quotient bar is missing in the OCR; replacing the quotient by the printed product changes the wave. The following M. Williams discussion of two interacting shocks is a separate item.

Use strain and momentum \(r_k=q_{k+1}-q_k\), \(p_k=\dot q_k\), which are localized for this wave; the displacement \(Q_\beta\) has nonzero limiting values and must not be treated as a localized strain. For a positive exponential weight define
\[
 \ell_a^2=\{v:\sum_{k\in\Z}e^{2ak}|v_k|^2<\infty\}.
\]
In the associated traveling-profile spectral calculation, a weighted far-field mode has exponent \(\nu=-a+i\xi\), \(\xi\in\R\). The aid gives the dispersion curves
\[
 \lambda_\pm(\xi)=c(-a+i\xi)\pm2\sinh\!\left(\frac{-a+i\xi}{2}\right),
\]
with spectral abscissa
\[
 b(a)=-ca+2\sinh(a/2).
\]
Here traveling-profile spectral parameters and discrete-time lattice Floquet multipliers must be related by the co-moving convention, rather than identified as the same spectrum.

For the one-site co-moving map take \((Sv)_k=v_{k+1}\), the linearized evolution \(\mathcal U(t,0)\), and
\[
 \mathcal M=S\mathcal U(1/c,0).
\]
The shift realigns the wave after it has moved one site. The corresponding far-field multipliers are \(\mu_\pm(\xi)=\exp(\lambda_\pm(\xi)/c)\), with maximum modulus
\[
 \rho_{\mathrm{ess}}(\mathcal M)
 =\exp\!\left(-a+\frac2c\sinh(a/2)\right).
\]
Changing the shift or weight convention changes these formulas.

The aid's sharp synthesis is the equivalence, for \(a>0\), of
\[
 0<a<2\beta,\qquad b(a)<0,\qquad
 \rho_{\mathrm{ess}}(\mathcal M)<1,
\]
and membership of the phase/speed tangent plane in the positively weighted strain-momentum space. The soliton and its tangents decay at rate \(2\beta\), with polynomial factors allowed for the speed derivative. At \(a=2\beta\) the tangent localization fails, the far-field spectral abscissa reaches zero, and the essential Floquet multipliers reach the unit circle; for \(a>2\beta\) the essential spectral radius exceeds one. At \(a=0\) the strict essential gap is also absent. Translation and speed variation produce neutral generalized modes, so asymptotic stability must treat their finite-dimensional plane separately rather than assert strict decay of every perturbation.

The supplied aid attributes the historical spectral question and weighted asymptotic stability, including the complete traveling-Floquet spectral description, to Mizumachi and Pego. Exact functional hypotheses, modulation statements, and the full spectral theorem are not supplied, so the displayed far-field calculation is not asserted to replace their discrete-spectrum analysis. The sharp essential threshold alone would not exclude additional unstable point spectrum.

For source preservation, the appended M. Williams setup is a different problem about \(u_t+\partial_xf(u)=0\) and \(u_t^\varepsilon+\partial_xf(u^\varepsilon)=\varepsilon\Delta u^\varepsilon\), with \(u\in\R^2\), \(x\in\R\), and a two-shock figure. Neither that absent figure nor a shock-interaction claim is part of the Toda target.
\subsection{Short English statement}
The recovered question concerns Pego's Toda soliton, not the appended shock problem. The aid reports its known weighted stability and identifies \(0<a<2\beta\) as the common strict threshold for essential spectral decay, localization of phase/speed tangents, and co-moving Floquet multipliers inside the unit disk.
\subsection{Sources}
\begin{itemize}
\item Local registry, permanent UnsolvedMath ID \(20002443\), exact public code \texttt{catalogue entry}; canonical location \texttt{aim-pdes-notes.json}, \texttt{notes[54]}.
\item AIM workshop \emph{Stability criteria for multi-dimensional waves and patterns}, Toda item attributed to R. Pego; supplied reference \url{https://aimath.org/WWN/multidimwaves/multidimwaves.pdf}. The verification field reports inspection of page 3 recovering the fraction bar and separating M. Williams's following item. No PDF was accessed here.
\item Mizumachi and Pego, attributed by the aid for weighted asymptotic stability and the complete traveling-Floquet spectrum. No title, year, journal, or theorem number is supplied. M. Williams is named only for the following two-shock setup.
\item Supplied research attempt 1: \texttt{partial\_result}, run and registry status \texttt{solved}, rigor label \texttt{proved}, low novelty confidence; classification \texttt{SOLVED-IN-LITERATURE}, created and updated 14 August 2026. The aid is machine generated, not peer reviewed, and requires independent verification before citation.
\end{itemize}
\end{document}
